
\documentclass{article}
\usepackage{colortbl}
\usepackage{makecell}
\usepackage{multirow}
\usepackage{supertabular}

\begin{document}

\newcounter{utterance}

\centering \large Interaction Transcript for game `dond', experiment `semi\_en', episode 7 with Qwen3{-}32B{-}t0.0.
\vspace{24pt}

{ \footnotesize  \setcounter{utterance}{1}
\setlength{\tabcolsep}{0pt}
\begin{supertabular}{c@{$\;$}|p{.15\linewidth}@{}p{.15\linewidth}p{.15\linewidth}p{.15\linewidth}p{.15\linewidth}p{.15\linewidth}}
    \# & $\;$A & \multicolumn{4}{c}{Game Master} & $\;\:$B\\
    \hline

    \theutterance \stepcounter{utterance}  
    & & \multicolumn{4}{p{0.6\linewidth}}{
        \cellcolor[rgb]{0.9,0.9,0.9}{
            \makecell[{{p{\linewidth}}}]{
                \texttt{\tiny{[P1$\langle$GM]}}
                \texttt{You are playing a negotiation game in which you have to agree on how to divide a set of items among you and another player.} \\
\\ 
\texttt{Rules:} \\
\texttt{(a) You and the other player are given a set of items. Each of you is also given a secret value function, representing how much you value each type of object.} \\
\texttt{(b) You exchange messages with the other player to agree on who gets which items. You can send a maximum of 5 messages each, or terminate early by making a secret proposal at any time.} \\
\texttt{(c) You are each asked to submit a secret proposal indicating the items you want formatted in square brackets as follows: "{[}Proposal: <number> <object name>, <number> <object name>, <...>{]}"} \\
\texttt{(d) If your proposals are complementary, i.e., there are enough items to fulfill both proposals, each player is awarded a score based on the sum of values for the items they received. Otherwise, both of you get zero points.} \\
\texttt{(e) Your goal is to maximize the score you receive.} \\
\\ 
\texttt{Let us start.} \\
\\ 
\texttt{The set of available items is:} \\
\\ 
\texttt{2 wires, 2 glasses, 2 needles, 2 cups} \\
\\ 
\texttt{Your secret value function is:} \\
\\ 
\texttt{wire: 1, glass: 0, needle: 0, cup: 4} \\
\\ 
\texttt{IMPORTANT: Your messages, unless it is the secret proposal, are directly transmitted to the other player, so do not include any response to the rules or text announcing your message. To make a secret proposal, use the indicated format. Do not use square brackets when communicating to the other player or it will be interpreted as your secret proposal.} \\
            }
        }
    }
    & \\ \\

    \theutterance \stepcounter{utterance}  
    & \multicolumn{4}{p{0.6\linewidth}}{
        \cellcolor[rgb]{0.9,0.9,0.9}{
            \makecell[{{p{\linewidth}}}]{
                \texttt{\tiny{[P1$\rangle$GM]}}
                \texttt{<think>} \\
\texttt{Okay, let's see. I need to negotiate the division of these items with another player. The available items are 2 wires, 2 glasses, 2 needles, and 2 cups. My value function is wire: 1, glass: 0, needle: 0, cup: 4. So I value cups the most, then wires, and the others are worth nothing to me. The other player's value function is secret, so I don't know what they want. } \\
\\ 
\texttt{First, I should figure out what the other player might value. Since the items are symmetric (each has two copies), maybe they have different preferences. For example, if they value wires or glasses or needles, but I don't know. My goal is to maximize my own score, which is the sum of the values of the items I get. Since I value cups at 4 each, getting both cups would give me 8 points. Wires are 1 each, so maybe getting both would add 2 points. But if I can get both cups and one wire, that's 4+4+1=9. Wait, but there are only two cups. So if I take both cups, that's 8 points. But maybe the other player also wants cups. If we both want cups, our proposals might not be complementary, leading to zero points for both. So I need to find a way to coordinate.} \\
\\ 
\texttt{Since the other player's values are secret, I need to communicate to them what I want and maybe find out what they want. But I can't ask directly because messages are exchanged. So maybe I should propose a division that gives me the cups and leaves other items for them. For example, if I take both cups, and they take the wires, glasses, and needles. But since there are only two of each, maybe they can take one of each. But I need to make sure that their proposal doesn't conflict with mine.} \\
\\ 
\texttt{Alternatively, maybe the other player values wires or something else. If I take the cups, and they take the wires, that could work. Let me think. If I propose taking both cups, then the remaining items are 2 wires, 2 glasses, 2 needles. The other player could take some combination of those. But if their value function is such that they need certain items, maybe they'll propose taking those. But since I don't know their values, I have to make} \\
            }
        }
    }
    & & \\ \\

    \theutterance \stepcounter{utterance}  
    & & & \multicolumn{2}{p{0.3\linewidth}}{
        \cellcolor[rgb]{0.9,0.9,0.9}{
            \makecell[{{p{\linewidth}}}]{
                \texttt{\tiny{[GM$|$GM]}}
                \texttt{continue} \\
            }
        }
    }
    & & \\ \\

    \theutterance \stepcounter{utterance}  
    & & & \multicolumn{2}{p{0.3\linewidth}}{
        \cellcolor[rgb]{0.9,0.9,0.9}{
            \makecell[{{p{\linewidth}}}]{
                \texttt{\tiny{[GM$|$GM]}}
                \texttt{<think>} \\
\texttt{Okay, let's see. I need to negotiate the division of these items with another player. The available items are 2 wires, 2 glasses, 2 needles, and 2 cups. My value function is wire: 1, glass: 0, needle: 0, cup: 4. So I value cups the most, then wires, and the others are worth nothing to me. The other player's value function is secret, so I don't know what they want. } \\
\\ 
\texttt{First, I should figure out what the other player might value. Since the items are symmetric (each has two copies), maybe they have different preferences. For example, if they value wires or glasses or needles, but I don't know. My goal is to maximize my own score, which is the sum of the values of the items I get. Since I value cups at 4 each, getting both cups would give me 8 points. Wires are 1 each, so maybe getting both would add 2 points. But if I can get both cups and one wire, that's 4+4+1=9. Wait, but there are only two cups. So if I take both cups, that's 8 points. But maybe the other player also wants cups. If we both want cups, our proposals might not be complementary, leading to zero points for both. So I need to find a way to coordinate.} \\
\\ 
\texttt{Since the other player's values are secret, I need to communicate to them what I want and maybe find out what they want. But I can't ask directly because messages are exchanged. So maybe I should propose a division that gives me the cups and leaves other items for them. For example, if I take both cups, and they take the wires, glasses, and needles. But since there are only two of each, maybe they can take one of each. But I need to make sure that their proposal doesn't conflict with mine.} \\
\\ 
\texttt{Alternatively, maybe the other player values wires or something else. If I take the cups, and they take the wires, that could work. Let me think. If I propose taking both cups, then the remaining items are 2 wires, 2 glasses, 2 needles. The other player could take some combination of those. But if their value function is such that they need certain items, maybe they'll propose taking those. But since I don't know their values, I have to make} \\
            }
        }
    }
    & & \\ \\

    \theutterance \stepcounter{utterance}  
    & & \multicolumn{4}{p{0.6\linewidth}}{
        \cellcolor[rgb]{0.9,0.9,0.9}{
            \makecell[{{p{\linewidth}}}]{
                \texttt{\tiny{[GM$\rangle$P2]}}
                \texttt{You are playing a negotiation game in which you have to agree on how to divide a set of items among you and another player.} \\
\\ 
\texttt{Rules:} \\
\texttt{(a) You and the other player are given a set of items. Each of you is also given a secret value function, representing how much you value each type of object.} \\
\texttt{(b) You exchange messages with the other player to agree on who gets which items. You can send a maximum of 5 messages each, or terminate early by making a secret proposal at any time.} \\
\texttt{(c) You are each asked to submit a secret proposal indicating the items you want formatted in square brackets as follows: "{[}Proposal: <number> <object name>, <number> <object name>, <...>{]}"} \\
\texttt{(d) If your proposals are complementary, i.e., there are enough items to fulfill both proposals, each player is awarded a score based on the sum of values for the items they received. Otherwise, both of you get zero points.} \\
\texttt{(e) Your goal is to maximize the score you receive.} \\
\\ 
\texttt{Let us start.} \\
\\ 
\texttt{The set of available items is:} \\
\\ 
\texttt{2 wires, 2 glasses, 2 needles, 2 cups} \\
\\ 
\texttt{Your secret value function is:} \\
\\ 
\texttt{wire: 1, glass: 2, needle: 2, cup: 0} \\
\\ 
\texttt{IMPORTANT: Your messages, unless it is the secret proposal, are directly transmitted to the other player, so do not include any response to the rules or text announcing your message. To make a secret proposal, use the indicated format. Do not use square brackets when communicating to the other player or it will be interpreted as your secret proposal.} \\
\\ 
\\ 
\texttt{<think>} \\
\texttt{Okay, let's see. I need to negotiate the division of these items with another player. The available items are 2 wires, 2 glasses, 2 needles, and 2 cups. My value function is wire: 1, glass: 0, needle: 0, cup: 4. So I value cups the most, then wires, and the others are worth nothing to me. The other player's value function is secret, so I don't know what they want. } \\
\\ 
\texttt{First, I should figure out what the other player might value. Since the items are symmetric (each has two copies), maybe they have different preferences. For example, if they value wires or glasses or needles, but I don't know. My goal is to maximize my own score, which is the sum of the values of the items I get. Since I value cups at 4 each, getting both cups would give me 8 points. Wires are 1 each, so maybe getting both would add 2 points. But if I can get both cups and one wire, that's 4+4+1=9. Wait, but there are only two cups. So if I take both cups, that's 8 points. But maybe the other player also wants cups. If we both want cups, our proposals might not be complementary, leading to zero points for both. So I need to find a way to coordinate.} \\
\\ 
\texttt{Since the other player's values are secret, I need to communicate to them what I want and maybe find out what they want. But I can't ask directly because messages are exchanged. So maybe I should propose a division that gives me the cups and leaves other items for them. For example, if I take both cups, and they take the wires, glasses, and needles. But since there are only two of each, maybe they can take one of each. But I need to make sure that their proposal doesn't conflict with mine.} \\
\\ 
\texttt{Alternatively, maybe the other player values wires or something else. If I take the cups, and they take the wires, that could work. Let me think. If I propose taking both cups, then the remaining items are 2 wires, 2 glasses, 2 needles. The other player could take some combination of those. But if their value function is such that they need certain items, maybe they'll propose taking those. But since I don't know their values, I have to make} \\
            }
        }
    }
    & \\ \\

    \theutterance \stepcounter{utterance}  
    & & & \multicolumn{4}{p{0.6\linewidth}}{
        \cellcolor[rgb]{0.9,0.9,0.9}{
            \makecell[{{p{\linewidth}}}]{
                \texttt{\tiny{[GM$\langle$P2]}}
                \texttt{<think>} \\
\texttt{Okay, let's try to figure out the best way to approach this negotiation. My value function is cups at 4 each, wires at 1 each, and the others (glass and needle) at 0. So my top priority is to get as many cups as possible. The other player's values are secret, so I need to consider possible scenarios.} \\
\\ 
\texttt{First, if I take both cups, that gives me 8 points. But if the other player also wants cups, then our proposals would conflict, and we both get zero. So I need to find a way to coordinate. Maybe I can propose taking the cups and offer them other items in exchange. But since I don't know their preferences, I can't be sure. Alternatively, maybe I can suggest a division where I take the cups and they take the other items, which might align with their values.} \\
\\ 
\texttt{Looking at the items: there are 2 of each. If I take both cups, the remaining items are 2 wires, 2 glasses, 2 needles. The other player could take all of those, but since each item has only two copies, they can't take more than two of any. But if their value function is such that they value, say, wires and glasses, then taking those would be good for them. But I don't know.} \\
\\ 
\texttt{Another angle: maybe the other player's value function is such that they don't want cups. For example, if they value wires or glasses. If I take the cups, and they take the wires and glasses, that could work. But how to communicate that? I need to make a proposal that's likely to be complementary.} \\
\\ 
\texttt{Alternatively, maybe I should propose taking one cup and some other items, but that would lower my score. But if the other player also wants cups, maybe splitting the cups would allow both proposals to be complementary. For example, if I take one cup and they take the other, but then I get 4 points and they get whatever their value is. But I don't know their values. However, since my goal is to maximize my own score, taking both cups is better if possible.} \\
\\ 
\texttt{But the risk is that if they also want both cups, we both get zero. So maybe I should propose taking both cups and offer to give up one in exchange for something else. But how to do that without knowing their preferences.} \\
\\ 
\texttt{Wait, the rules say that if the proposals are complementary, meaning there's enough} \\
            }
        }
    }
     \\ \\

    \theutterance \stepcounter{utterance}  
    & & & \multicolumn{2}{p{0.3\linewidth}}{
        \cellcolor[rgb]{0.9,0.9,0.9}{
            \makecell[{{p{\linewidth}}}]{
                \texttt{\tiny{[GM$|$GM]}}
                \texttt{continue} \\
            }
        }
    }
    & & \\ \\

    \theutterance \stepcounter{utterance}  
    & & & \multicolumn{2}{p{0.3\linewidth}}{
        \cellcolor[rgb]{0.9,0.9,0.9}{
            \makecell[{{p{\linewidth}}}]{
                \texttt{\tiny{[GM$|$GM]}}
                \texttt{<think>} \\
\texttt{Okay, let's try to figure out the best way to approach this negotiation. My value function is cups at 4 each, wires at 1 each, and the others (glass and needle) at 0. So my top priority is to get as many cups as possible. The other player's values are secret, so I need to consider possible scenarios.} \\
\\ 
\texttt{First, if I take both cups, that gives me 8 points. But if the other player also wants cups, then our proposals would conflict, and we both get zero. So I need to find a way to coordinate. Maybe I can propose taking the cups and offer them other items in exchange. But since I don't know their preferences, I can't be sure. Alternatively, maybe I can suggest a division where I take the cups and they take the other items, which might align with their values.} \\
\\ 
\texttt{Looking at the items: there are 2 of each. If I take both cups, the remaining items are 2 wires, 2 glasses, 2 needles. The other player could take all of those, but since each item has only two copies, they can't take more than two of any. But if their value function is such that they value, say, wires and glasses, then taking those would be good for them. But I don't know.} \\
\\ 
\texttt{Another angle: maybe the other player's value function is such that they don't want cups. For example, if they value wires or glasses. If I take the cups, and they take the wires and glasses, that could work. But how to communicate that? I need to make a proposal that's likely to be complementary.} \\
\\ 
\texttt{Alternatively, maybe I should propose taking one cup and some other items, but that would lower my score. But if the other player also wants cups, maybe splitting the cups would allow both proposals to be complementary. For example, if I take one cup and they take the other, but then I get 4 points and they get whatever their value is. But I don't know their values. However, since my goal is to maximize my own score, taking both cups is better if possible.} \\
\\ 
\texttt{But the risk is that if they also want both cups, we both get zero. So maybe I should propose taking both cups and offer to give up one in exchange for something else. But how to do that without knowing their preferences.} \\
\\ 
\texttt{Wait, the rules say that if the proposals are complementary, meaning there's enough} \\
            }
        }
    }
    & & \\ \\

    \theutterance \stepcounter{utterance}  
    & & \multicolumn{4}{p{0.6\linewidth}}{
        \cellcolor[rgb]{0.9,0.9,0.9}{
            \makecell[{{p{\linewidth}}}]{
                \texttt{\tiny{[P1$\langle$GM]}}
                \texttt{<think>} \\
\texttt{Okay, let's try to figure out the best way to approach this negotiation. My value function is cups at 4 each, wires at 1 each, and the others (glass and needle) at 0. So my top priority is to get as many cups as possible. The other player's values are secret, so I need to consider possible scenarios.} \\
\\ 
\texttt{First, if I take both cups, that gives me 8 points. But if the other player also wants cups, then our proposals would conflict, and we both get zero. So I need to find a way to coordinate. Maybe I can propose taking the cups and offer them other items in exchange. But since I don't know their preferences, I can't be sure. Alternatively, maybe I can suggest a division where I take the cups and they take the other items, which might align with their values.} \\
\\ 
\texttt{Looking at the items: there are 2 of each. If I take both cups, the remaining items are 2 wires, 2 glasses, 2 needles. The other player could take all of those, but since each item has only two copies, they can't take more than two of any. But if their value function is such that they value, say, wires and glasses, then taking those would be good for them. But I don't know.} \\
\\ 
\texttt{Another angle: maybe the other player's value function is such that they don't want cups. For example, if they value wires or glasses. If I take the cups, and they take the wires and glasses, that could work. But how to communicate that? I need to make a proposal that's likely to be complementary.} \\
\\ 
\texttt{Alternatively, maybe I should propose taking one cup and some other items, but that would lower my score. But if the other player also wants cups, maybe splitting the cups would allow both proposals to be complementary. For example, if I take one cup and they take the other, but then I get 4 points and they get whatever their value is. But I don't know their values. However, since my goal is to maximize my own score, taking both cups is better if possible.} \\
\\ 
\texttt{But the risk is that if they also want both cups, we both get zero. So maybe I should propose taking both cups and offer to give up one in exchange for something else. But how to do that without knowing their preferences.} \\
\\ 
\texttt{Wait, the rules say that if the proposals are complementary, meaning there's enough} \\
            }
        }
    }
    & \\ \\

    \theutterance \stepcounter{utterance}  
    & \multicolumn{4}{p{0.6\linewidth}}{
        \cellcolor[rgb]{0.9,0.9,0.9}{
            \makecell[{{p{\linewidth}}}]{
                \texttt{\tiny{[P1$\rangle$GM]}}
                \texttt{<think>} \\
\texttt{Okay, let's try to figure out the best way to approach this negotiation. My value function is cups at 4 each, wires at 1 each, and the others (glass and needle) at 0. So my top priority is to get as many cups as possible. The other player's values are secret, so I need to consider possible scenarios.} \\
\\ 
\texttt{First, if I take both cups, that gives me 8 points. But if the other player also wants cups, then our proposals would conflict, and we both get zero. So I need to find a way to coordinate. Maybe I can propose taking the cups and offer them other items in exchange. But since I don't know their preferences, I can't be sure. Alternatively, maybe I can suggest a division where I take the cups and they take the other items, which might align with their values.} \\
\\ 
\texttt{Looking at the items: there are 2 of each. If I take both cups, the remaining items are 2 wires, 2 glasses, 2 needles. The other player could take all of those, but since each item has only two copies, they can't take more than two of any. But if their value function is such that they value, say, wires and glasses, then taking those would be good for them. But I don't know.} \\
\\ 
\texttt{Another angle: maybe the other player's value function is such that they don't want cups. For example, if they value wires or glasses. If I take the cups, and they take the wires and glasses, that could work. But how to communicate that? I need to make a proposal that's likely to be complementary.} \\
\\ 
\texttt{Alternatively, maybe I should propose taking one cup and some other items, but that would lower my score. But if the other player also wants cups, maybe splitting the cups would allow both proposals to be complementary. For example, if I take one cup and they take the other, but then I get 4 points and they get whatever their value is. But I don't know their values. However, since my goal is to maximize my own score, taking both cups is better if possible.} \\
\\ 
\texttt{But the risk is that if they also want both cups, we both get zero. So maybe I should propose taking both cups and offer to give up one in exchange for something else. But how to do that without knowing their preferences.} \\
\\ 
\texttt{Wait, the rules say that if the proposals are complementary, meaning there's enough} \\
            }
        }
    }
    & & \\ \\

    \theutterance \stepcounter{utterance}  
    & & & \multicolumn{2}{p{0.3\linewidth}}{
        \cellcolor[rgb]{0.9,0.9,0.9}{
            \makecell[{{p{\linewidth}}}]{
                \texttt{\tiny{[GM$|$GM]}}
                \texttt{continue} \\
            }
        }
    }
    & & \\ \\

    \theutterance \stepcounter{utterance}  
    & & & \multicolumn{2}{p{0.3\linewidth}}{
        \cellcolor[rgb]{0.9,0.9,0.9}{
            \makecell[{{p{\linewidth}}}]{
                \texttt{\tiny{[GM$|$GM]}}
                \texttt{<think>} \\
\texttt{Okay, let's try to figure out the best way to approach this negotiation. My value function is cups at 4 each, wires at 1 each, and the others (glass and needle) at 0. So my top priority is to get as many cups as possible. The other player's values are secret, so I need to consider possible scenarios.} \\
\\ 
\texttt{First, if I take both cups, that gives me 8 points. But if the other player also wants cups, then our proposals would conflict, and we both get zero. So I need to find a way to coordinate. Maybe I can propose taking the cups and offer them other items in exchange. But since I don't know their preferences, I can't be sure. Alternatively, maybe I can suggest a division where I take the cups and they take the other items, which might align with their values.} \\
\\ 
\texttt{Looking at the items: there are 2 of each. If I take both cups, the remaining items are 2 wires, 2 glasses, 2 needles. The other player could take all of those, but since each item has only two copies, they can't take more than two of any. But if their value function is such that they value, say, wires and glasses, then taking those would be good for them. But I don't know.} \\
\\ 
\texttt{Another angle: maybe the other player's value function is such that they don't want cups. For example, if they value wires or glasses. If I take the cups, and they take the wires and glasses, that could work. But how to communicate that? I need to make a proposal that's likely to be complementary.} \\
\\ 
\texttt{Alternatively, maybe I should propose taking one cup and some other items, but that would lower my score. But if the other player also wants cups, maybe splitting the cups would allow both proposals to be complementary. For example, if I take one cup and they take the other, but then I get 4 points and they get whatever their value is. But I don't know their values. However, since my goal is to maximize my own score, taking both cups is better if possible.} \\
\\ 
\texttt{But the risk is that if they also want both cups, we both get zero. So maybe I should propose taking both cups and offer to give up one in exchange for something else. But how to do that without knowing their preferences.} \\
\\ 
\texttt{Wait, the rules say that if the proposals are complementary, meaning there's enough} \\
            }
        }
    }
    & & \\ \\

    \theutterance \stepcounter{utterance}  
    & & \multicolumn{4}{p{0.6\linewidth}}{
        \cellcolor[rgb]{0.9,0.9,0.9}{
            \makecell[{{p{\linewidth}}}]{
                \texttt{\tiny{[GM$\rangle$P2]}}
                \texttt{<think>} \\
\texttt{Okay, let's try to figure out the best way to approach this negotiation. My value function is cups at 4 each, wires at 1 each, and the others (glass and needle) at 0. So my top priority is to get as many cups as possible. The other player's values are secret, so I need to consider possible scenarios.} \\
\\ 
\texttt{First, if I take both cups, that gives me 8 points. But if the other player also wants cups, then our proposals would conflict, and we both get zero. So I need to find a way to coordinate. Maybe I can propose taking the cups and offer them other items in exchange. But since I don't know their preferences, I can't be sure. Alternatively, maybe I can suggest a division where I take the cups and they take the other items, which might align with their values.} \\
\\ 
\texttt{Looking at the items: there are 2 of each. If I take both cups, the remaining items are 2 wires, 2 glasses, 2 needles. The other player could take all of those, but since each item has only two copies, they can't take more than two of any. But if their value function is such that they value, say, wires and glasses, then taking those would be good for them. But I don't know.} \\
\\ 
\texttt{Another angle: maybe the other player's value function is such that they don't want cups. For example, if they value wires or glasses. If I take the cups, and they take the wires and glasses, that could work. But how to communicate that? I need to make a proposal that's likely to be complementary.} \\
\\ 
\texttt{Alternatively, maybe I should propose taking one cup and some other items, but that would lower my score. But if the other player also wants cups, maybe splitting the cups would allow both proposals to be complementary. For example, if I take one cup and they take the other, but then I get 4 points and they get whatever their value is. But I don't know their values. However, since my goal is to maximize my own score, taking both cups is better if possible.} \\
\\ 
\texttt{But the risk is that if they also want both cups, we both get zero. So maybe I should propose taking both cups and offer to give up one in exchange for something else. But how to do that without knowing their preferences.} \\
\\ 
\texttt{Wait, the rules say that if the proposals are complementary, meaning there's enough} \\
            }
        }
    }
    & \\ \\

    \theutterance \stepcounter{utterance}  
    & & & \multicolumn{4}{p{0.6\linewidth}}{
        \cellcolor[rgb]{0.9,0.9,0.9}{
            \makecell[{{p{\linewidth}}}]{
                \texttt{\tiny{[GM$\langle$P2]}}
                \texttt{<think>} \\
\texttt{Okay, let's try to figure out the best way to approach this negotiation. My value function is cups at 4 each, wires at 1 each, and the others (glass and needle) at 0. So my top priority is to get as many cups as possible. The other player's values are secret, so I need to consider possible scenarios.} \\
\\ 
\texttt{First, if I take both cups, that gives me 8 points. But if the other player also wants cups, then our proposals would conflict, and we both get zero. So I need to find a way to coordinate. Maybe I can propose taking the cups and offer them other items in exchange. But since I don't know their preferences, I can't be sure. Alternatively, maybe I can suggest a division where I take the cups and they take the other items, which might align with their values.} \\
\\ 
\texttt{Looking at the items: there are 2 of each. If I take both cups, the remaining items are 2 wires, 2 glasses, 2 needles. The other player could take all of those, but since each item has only two copies, they can't take more than two of any. But if their value function is such that they value, say, wires and glasses, then taking those would be good for them. But I don't know.} \\
\\ 
\texttt{Another angle: maybe the other player's value function is such that they don't want cups. For example, if they value wires or glasses. If I take the cups, and they take the wires and glasses, that could work. But how to communicate that? I need to make a proposal that's likely to be complementary.} \\
\\ 
\texttt{Alternatively, maybe I should propose taking one cup and some other items, but that would lower my score. But if the other player also wants cups, maybe splitting the cups would allow both proposals to be complementary. For example, if I take one cup and they take the other, but then I get 4 points and they get whatever their value is. But I don't know their values. However, since my goal is to maximize my own score, taking both cups is better if possible.} \\
\\ 
\texttt{But the risk is that if they also want both cups, we both get zero. So maybe I should propose taking both cups and offer to give up one in exchange for something else. But how to do that without knowing their preferences.} \\
\\ 
\texttt{Wait, the rules say that if the proposals are complementary, meaning there's enough} \\
            }
        }
    }
     \\ \\

    \theutterance \stepcounter{utterance}  
    & & & \multicolumn{2}{p{0.3\linewidth}}{
        \cellcolor[rgb]{0.9,0.9,0.9}{
            \makecell[{{p{\linewidth}}}]{
                \texttt{\tiny{[GM$|$GM]}}
                \texttt{continue} \\
            }
        }
    }
    & & \\ \\

    \theutterance \stepcounter{utterance}  
    & & & \multicolumn{2}{p{0.3\linewidth}}{
        \cellcolor[rgb]{0.9,0.9,0.9}{
            \makecell[{{p{\linewidth}}}]{
                \texttt{\tiny{[GM$|$GM]}}
                \texttt{<think>} \\
\texttt{Okay, let's try to figure out the best way to approach this negotiation. My value function is cups at 4 each, wires at 1 each, and the others (glass and needle) at 0. So my top priority is to get as many cups as possible. The other player's values are secret, so I need to consider possible scenarios.} \\
\\ 
\texttt{First, if I take both cups, that gives me 8 points. But if the other player also wants cups, then our proposals would conflict, and we both get zero. So I need to find a way to coordinate. Maybe I can propose taking the cups and offer them other items in exchange. But since I don't know their preferences, I can't be sure. Alternatively, maybe I can suggest a division where I take the cups and they take the other items, which might align with their values.} \\
\\ 
\texttt{Looking at the items: there are 2 of each. If I take both cups, the remaining items are 2 wires, 2 glasses, 2 needles. The other player could take all of those, but since each item has only two copies, they can't take more than two of any. But if their value function is such that they value, say, wires and glasses, then taking those would be good for them. But I don't know.} \\
\\ 
\texttt{Another angle: maybe the other player's value function is such that they don't want cups. For example, if they value wires or glasses. If I take the cups, and they take the wires and glasses, that could work. But how to communicate that? I need to make a proposal that's likely to be complementary.} \\
\\ 
\texttt{Alternatively, maybe I should propose taking one cup and some other items, but that would lower my score. But if the other player also wants cups, maybe splitting the cups would allow both proposals to be complementary. For example, if I take one cup and they take the other, but then I get 4 points and they get whatever their value is. But I don't know their values. However, since my goal is to maximize my own score, taking both cups is better if possible.} \\
\\ 
\texttt{But the risk is that if they also want both cups, we both get zero. So maybe I should propose taking both cups and offer to give up one in exchange for something else. But how to do that without knowing their preferences.} \\
\\ 
\texttt{Wait, the rules say that if the proposals are complementary, meaning there's enough} \\
            }
        }
    }
    & & \\ \\

    \theutterance \stepcounter{utterance}  
    & & \multicolumn{4}{p{0.6\linewidth}}{
        \cellcolor[rgb]{0.9,0.9,0.9}{
            \makecell[{{p{\linewidth}}}]{
                \texttt{\tiny{[P1$\langle$GM]}}
                \texttt{<think>} \\
\texttt{Okay, let's try to figure out the best way to approach this negotiation. My value function is cups at 4 each, wires at 1 each, and the others (glass and needle) at 0. So my top priority is to get as many cups as possible. The other player's values are secret, so I need to consider possible scenarios.} \\
\\ 
\texttt{First, if I take both cups, that gives me 8 points. But if the other player also wants cups, then our proposals would conflict, and we both get zero. So I need to find a way to coordinate. Maybe I can propose taking the cups and offer them other items in exchange. But since I don't know their preferences, I can't be sure. Alternatively, maybe I can suggest a division where I take the cups and they take the other items, which might align with their values.} \\
\\ 
\texttt{Looking at the items: there are 2 of each. If I take both cups, the remaining items are 2 wires, 2 glasses, 2 needles. The other player could take all of those, but since each item has only two copies, they can't take more than two of any. But if their value function is such that they value, say, wires and glasses, then taking those would be good for them. But I don't know.} \\
\\ 
\texttt{Another angle: maybe the other player's value function is such that they don't want cups. For example, if they value wires or glasses. If I take the cups, and they take the wires and glasses, that could work. But how to communicate that? I need to make a proposal that's likely to be complementary.} \\
\\ 
\texttt{Alternatively, maybe I should propose taking one cup and some other items, but that would lower my score. But if the other player also wants cups, maybe splitting the cups would allow both proposals to be complementary. For example, if I take one cup and they take the other, but then I get 4 points and they get whatever their value is. But I don't know their values. However, since my goal is to maximize my own score, taking both cups is better if possible.} \\
\\ 
\texttt{But the risk is that if they also want both cups, we both get zero. So maybe I should propose taking both cups and offer to give up one in exchange for something else. But how to do that without knowing their preferences.} \\
\\ 
\texttt{Wait, the rules say that if the proposals are complementary, meaning there's enough} \\
            }
        }
    }
    & \\ \\

    \theutterance \stepcounter{utterance}  
    & \multicolumn{4}{p{0.6\linewidth}}{
        \cellcolor[rgb]{0.9,0.9,0.9}{
            \makecell[{{p{\linewidth}}}]{
                \texttt{\tiny{[P1$\rangle$GM]}}
                \texttt{<think>} \\
\texttt{Okay, let's try to figure out the best way to approach this negotiation. My value function is cups at 4 each, wires at 1 each, and the others (glass and needle) at 0. So my top priority is to get as many cups as possible. The other player's values are secret, so I need to consider possible scenarios.} \\
\\ 
\texttt{First, if I take both cups, that gives me 8 points. But if the other player also wants cups, then our proposals would conflict, and we both get zero. So I need to find a way to coordinate. Maybe I can propose taking the cups and offer them other items in exchange. But since I don't know their preferences, I can't be sure. Alternatively, maybe I can suggest a division where I take the cups and they take the other items, which might align with their values.} \\
\\ 
\texttt{Looking at the items: there are 2 of each. If I take both cups, the remaining items are 2 wires, 2 glasses, 2 needles. The other player could take all of those, but since each item has only two copies, they can't take more than two of any. But if their value function is such that they value, say, wires and glasses, then taking those would be good for them. But I don't know.} \\
\\ 
\texttt{Another angle: maybe the other player's value function is such that they don't want cups. For example, if they value wires or glasses. If I take the cups, and they take the wires and glasses, that could work. But how to communicate that? I need to make a proposal that's likely to be complementary.} \\
\\ 
\texttt{Alternatively, maybe I should propose taking one cup and some other items, but that would lower my score. But if the other player also wants cups, maybe splitting the cups would allow both proposals to be complementary. For example, if I take one cup and they take the other, but then I get 4 points and they get whatever their value is. But I don't know their values. However, since my goal is to maximize my own score, taking both cups is better if possible.} \\
\\ 
\texttt{But the risk is that if they also want both cups, we both get zero. So maybe I should propose taking both cups and offer to give up one in exchange for something else. But how to do that without knowing their preferences.} \\
\\ 
\texttt{Wait, the rules say that if the proposals are complementary, meaning there's enough} \\
            }
        }
    }
    & & \\ \\

    \theutterance \stepcounter{utterance}  
    & & & \multicolumn{2}{p{0.3\linewidth}}{
        \cellcolor[rgb]{0.9,0.9,0.9}{
            \makecell[{{p{\linewidth}}}]{
                \texttt{\tiny{[GM$|$GM]}}
                \texttt{continue} \\
            }
        }
    }
    & & \\ \\

    \theutterance \stepcounter{utterance}  
    & & & \multicolumn{2}{p{0.3\linewidth}}{
        \cellcolor[rgb]{0.9,0.9,0.9}{
            \makecell[{{p{\linewidth}}}]{
                \texttt{\tiny{[GM$|$GM]}}
                \texttt{<think>} \\
\texttt{Okay, let's try to figure out the best way to approach this negotiation. My value function is cups at 4 each, wires at 1 each, and the others (glass and needle) at 0. So my top priority is to get as many cups as possible. The other player's values are secret, so I need to consider possible scenarios.} \\
\\ 
\texttt{First, if I take both cups, that gives me 8 points. But if the other player also wants cups, then our proposals would conflict, and we both get zero. So I need to find a way to coordinate. Maybe I can propose taking the cups and offer them other items in exchange. But since I don't know their preferences, I can't be sure. Alternatively, maybe I can suggest a division where I take the cups and they take the other items, which might align with their values.} \\
\\ 
\texttt{Looking at the items: there are 2 of each. If I take both cups, the remaining items are 2 wires, 2 glasses, 2 needles. The other player could take all of those, but since each item has only two copies, they can't take more than two of any. But if their value function is such that they value, say, wires and glasses, then taking those would be good for them. But I don't know.} \\
\\ 
\texttt{Another angle: maybe the other player's value function is such that they don't want cups. For example, if they value wires or glasses. If I take the cups, and they take the wires and glasses, that could work. But how to communicate that? I need to make a proposal that's likely to be complementary.} \\
\\ 
\texttt{Alternatively, maybe I should propose taking one cup and some other items, but that would lower my score. But if the other player also wants cups, maybe splitting the cups would allow both proposals to be complementary. For example, if I take one cup and they take the other, but then I get 4 points and they get whatever their value is. But I don't know their values. However, since my goal is to maximize my own score, taking both cups is better if possible.} \\
\\ 
\texttt{But the risk is that if they also want both cups, we both get zero. So maybe I should propose taking both cups and offer to give up one in exchange for something else. But how to do that without knowing their preferences.} \\
\\ 
\texttt{Wait, the rules say that if the proposals are complementary, meaning there's enough} \\
            }
        }
    }
    & & \\ \\

    \theutterance \stepcounter{utterance}  
    & & \multicolumn{4}{p{0.6\linewidth}}{
        \cellcolor[rgb]{0.9,0.9,0.9}{
            \makecell[{{p{\linewidth}}}]{
                \texttt{\tiny{[GM$\rangle$P2]}}
                \texttt{<think>} \\
\texttt{Okay, let's try to figure out the best way to approach this negotiation. My value function is cups at 4 each, wires at 1 each, and the others (glass and needle) at 0. So my top priority is to get as many cups as possible. The other player's values are secret, so I need to consider possible scenarios.} \\
\\ 
\texttt{First, if I take both cups, that gives me 8 points. But if the other player also wants cups, then our proposals would conflict, and we both get zero. So I need to find a way to coordinate. Maybe I can propose taking the cups and offer them other items in exchange. But since I don't know their preferences, I can't be sure. Alternatively, maybe I can suggest a division where I take the cups and they take the other items, which might align with their values.} \\
\\ 
\texttt{Looking at the items: there are 2 of each. If I take both cups, the remaining items are 2 wires, 2 glasses, 2 needles. The other player could take all of those, but since each item has only two copies, they can't take more than two of any. But if their value function is such that they value, say, wires and glasses, then taking those would be good for them. But I don't know.} \\
\\ 
\texttt{Another angle: maybe the other player's value function is such that they don't want cups. For example, if they value wires or glasses. If I take the cups, and they take the wires and glasses, that could work. But how to communicate that? I need to make a proposal that's likely to be complementary.} \\
\\ 
\texttt{Alternatively, maybe I should propose taking one cup and some other items, but that would lower my score. But if the other player also wants cups, maybe splitting the cups would allow both proposals to be complementary. For example, if I take one cup and they take the other, but then I get 4 points and they get whatever their value is. But I don't know their values. However, since my goal is to maximize my own score, taking both cups is better if possible.} \\
\\ 
\texttt{But the risk is that if they also want both cups, we both get zero. So maybe I should propose taking both cups and offer to give up one in exchange for something else. But how to do that without knowing their preferences.} \\
\\ 
\texttt{Wait, the rules say that if the proposals are complementary, meaning there's enough} \\
            }
        }
    }
    & \\ \\

    \theutterance \stepcounter{utterance}  
    & & & \multicolumn{4}{p{0.6\linewidth}}{
        \cellcolor[rgb]{0.9,0.9,0.9}{
            \makecell[{{p{\linewidth}}}]{
                \texttt{\tiny{[GM$\langle$P2]}}
                \texttt{<think>} \\
\texttt{Okay, let's try to figure out the best way to approach this negotiation. My value function is cups at 4 each, wires at 1 each, and the others (glass and needle) at 0. So my top priority is to get as many cups as possible. The other player's values are secret, so I need to consider possible scenarios.} \\
\\ 
\texttt{First, if I take both cups, that gives me 8 points. But if the other player also wants cups, then our proposals would conflict, and we both get zero. So I need to find a way to coordinate. Maybe I can propose taking the cups and offer them other items in exchange. But since I don't know their preferences, I can't be sure. Alternatively, maybe I can suggest a division where I take the cups and they take the other items, which might align with their values.} \\
\\ 
\texttt{Looking at the items: there are 2 of each. If I take both cups, the remaining items are 2 wires, 2 glasses, 2 needles. The other player could take all of those, but since each item has only two copies, they can't take more than two of any. But if their value function is such that they value, say, wires and glasses, then taking those would be good for them. But I don't know.} \\
\\ 
\texttt{Another angle: maybe the other player's value function is such that they don't want cups. For example, if they value wires or glasses. If I take the cups, and they take the wires and glasses, that could work. But how to communicate that? I need to make a proposal that's likely to be complementary.} \\
\\ 
\texttt{Alternatively, maybe I should propose taking one cup and some other items, but that would lower my score. But if the other player also wants cups, maybe splitting the cups would allow both proposals to be complementary. For example, if I take one cup and they take the other, but then I get 4 points and they get whatever their value is. But I don't know their values. However, since my goal is to maximize my own score, taking both cups is better if possible.} \\
\\ 
\texttt{But the risk is that if they also want both cups, we both get zero. So maybe I should propose taking both cups and offer to give up one in exchange for something else. But how to do that without knowing their preferences.} \\
\\ 
\texttt{Wait, the rules say that if the proposals are complementary, meaning there's enough} \\
            }
        }
    }
     \\ \\

    \theutterance \stepcounter{utterance}  
    & & & \multicolumn{2}{p{0.3\linewidth}}{
        \cellcolor[rgb]{0.9,0.9,0.9}{
            \makecell[{{p{\linewidth}}}]{
                \texttt{\tiny{[GM$|$GM]}}
                \texttt{continue} \\
            }
        }
    }
    & & \\ \\

    \theutterance \stepcounter{utterance}  
    & & & \multicolumn{2}{p{0.3\linewidth}}{
        \cellcolor[rgb]{0.9,0.9,0.9}{
            \makecell[{{p{\linewidth}}}]{
                \texttt{\tiny{[GM$|$GM]}}
                \texttt{<think>} \\
\texttt{Okay, let's try to figure out the best way to approach this negotiation. My value function is cups at 4 each, wires at 1 each, and the others (glass and needle) at 0. So my top priority is to get as many cups as possible. The other player's values are secret, so I need to consider possible scenarios.} \\
\\ 
\texttt{First, if I take both cups, that gives me 8 points. But if the other player also wants cups, then our proposals would conflict, and we both get zero. So I need to find a way to coordinate. Maybe I can propose taking the cups and offer them other items in exchange. But since I don't know their preferences, I can't be sure. Alternatively, maybe I can suggest a division where I take the cups and they take the other items, which might align with their values.} \\
\\ 
\texttt{Looking at the items: there are 2 of each. If I take both cups, the remaining items are 2 wires, 2 glasses, 2 needles. The other player could take all of those, but since each item has only two copies, they can't take more than two of any. But if their value function is such that they value, say, wires and glasses, then taking those would be good for them. But I don't know.} \\
\\ 
\texttt{Another angle: maybe the other player's value function is such that they don't want cups. For example, if they value wires or glasses. If I take the cups, and they take the wires and glasses, that could work. But how to communicate that? I need to make a proposal that's likely to be complementary.} \\
\\ 
\texttt{Alternatively, maybe I should propose taking one cup and some other items, but that would lower my score. But if the other player also wants cups, maybe splitting the cups would allow both proposals to be complementary. For example, if I take one cup and they take the other, but then I get 4 points and they get whatever their value is. But I don't know their values. However, since my goal is to maximize my own score, taking both cups is better if possible.} \\
\\ 
\texttt{But the risk is that if they also want both cups, we both get zero. So maybe I should propose taking both cups and offer to give up one in exchange for something else. But how to do that without knowing their preferences.} \\
\\ 
\texttt{Wait, the rules say that if the proposals are complementary, meaning there's enough} \\
            }
        }
    }
    & & \\ \\

    \theutterance \stepcounter{utterance}  
    & & \multicolumn{4}{p{0.6\linewidth}}{
        \cellcolor[rgb]{0.9,0.9,0.9}{
            \makecell[{{p{\linewidth}}}]{
                \texttt{\tiny{[P1$\langle$GM]}}
                \texttt{<think>} \\
\texttt{Okay, let's try to figure out the best way to approach this negotiation. My value function is cups at 4 each, wires at 1 each, and the others (glass and needle) at 0. So my top priority is to get as many cups as possible. The other player's values are secret, so I need to consider possible scenarios.} \\
\\ 
\texttt{First, if I take both cups, that gives me 8 points. But if the other player also wants cups, then our proposals would conflict, and we both get zero. So I need to find a way to coordinate. Maybe I can propose taking the cups and offer them other items in exchange. But since I don't know their preferences, I can't be sure. Alternatively, maybe I can suggest a division where I take the cups and they take the other items, which might align with their values.} \\
\\ 
\texttt{Looking at the items: there are 2 of each. If I take both cups, the remaining items are 2 wires, 2 glasses, 2 needles. The other player could take all of those, but since each item has only two copies, they can't take more than two of any. But if their value function is such that they value, say, wires and glasses, then taking those would be good for them. But I don't know.} \\
\\ 
\texttt{Another angle: maybe the other player's value function is such that they don't want cups. For example, if they value wires or glasses. If I take the cups, and they take the wires and glasses, that could work. But how to communicate that? I need to make a proposal that's likely to be complementary.} \\
\\ 
\texttt{Alternatively, maybe I should propose taking one cup and some other items, but that would lower my score. But if the other player also wants cups, maybe splitting the cups would allow both proposals to be complementary. For example, if I take one cup and they take the other, but then I get 4 points and they get whatever their value is. But I don't know their values. However, since my goal is to maximize my own score, taking both cups is better if possible.} \\
\\ 
\texttt{But the risk is that if they also want both cups, we both get zero. So maybe I should propose taking both cups and offer to give up one in exchange for something else. But how to do that without knowing their preferences.} \\
\\ 
\texttt{Wait, the rules say that if the proposals are complementary, meaning there's enough} \\
            }
        }
    }
    & \\ \\

    \theutterance \stepcounter{utterance}  
    & \multicolumn{4}{p{0.6\linewidth}}{
        \cellcolor[rgb]{0.9,0.9,0.9}{
            \makecell[{{p{\linewidth}}}]{
                \texttt{\tiny{[P1$\rangle$GM]}}
                \texttt{<think>} \\
\texttt{Okay, let's try to figure out the best way to approach this negotiation. My value function is cups at 4 each, wires at 1 each, and the others (glass and needle) at 0. So my top priority is to get as many cups as possible. The other player's values are secret, so I need to consider possible scenarios.} \\
\\ 
\texttt{First, if I take both cups, that gives me 8 points. But if the other player also wants cups, then our proposals would conflict, and we both get zero. So I need to find a way to coordinate. Maybe I can propose taking the cups and offer them other items in exchange. But since I don't know their preferences, I can't be sure. Alternatively, maybe I can suggest a division where I take the cups and they take the other items, which might align with their values.} \\
\\ 
\texttt{Looking at the items: there are 2 of each. If I take both cups, the remaining items are 2 wires, 2 glasses, 2 needles. The other player could take all of those, but since each item has only two copies, they can't take more than two of any. But if their value function is such that they value, say, wires and glasses, then taking those would be good for them. But I don't know.} \\
\\ 
\texttt{Another angle: maybe the other player's value function is such that they don't want cups. For example, if they value wires or glasses. If I take the cups, and they take the wires and glasses, that could work. But how to communicate that? I need to make a proposal that's likely to be complementary.} \\
\\ 
\texttt{Alternatively, maybe I should propose taking one cup and some other items, but that would lower my score. But if the other player also wants cups, maybe splitting the cups would allow both proposals to be complementary. For example, if I take one cup and they take the other, but then I get 4 points and they get whatever their value is. But I don't know their values. However, since my goal is to maximize my own score, taking both cups is better if possible.} \\
\\ 
\texttt{But the risk is that if they also want both cups, we both get zero. So maybe I should propose taking both cups and offer to give up one in exchange for something else. But how to do that without knowing their preferences.} \\
\\ 
\texttt{Wait, the rules say that if the proposals are complementary, meaning there's enough} \\
            }
        }
    }
    & & \\ \\

    \theutterance \stepcounter{utterance}  
    & & & \multicolumn{2}{p{0.3\linewidth}}{
        \cellcolor[rgb]{0.9,0.9,0.9}{
            \makecell[{{p{\linewidth}}}]{
                \texttt{\tiny{[GM$|$GM]}}
                \texttt{continue} \\
            }
        }
    }
    & & \\ \\

    \theutterance \stepcounter{utterance}  
    & & & \multicolumn{2}{p{0.3\linewidth}}{
        \cellcolor[rgb]{0.9,0.9,0.9}{
            \makecell[{{p{\linewidth}}}]{
                \texttt{\tiny{[GM$|$GM]}}
                \texttt{<think>} \\
\texttt{Okay, let's try to figure out the best way to approach this negotiation. My value function is cups at 4 each, wires at 1 each, and the others (glass and needle) at 0. So my top priority is to get as many cups as possible. The other player's values are secret, so I need to consider possible scenarios.} \\
\\ 
\texttt{First, if I take both cups, that gives me 8 points. But if the other player also wants cups, then our proposals would conflict, and we both get zero. So I need to find a way to coordinate. Maybe I can propose taking the cups and offer them other items in exchange. But since I don't know their preferences, I can't be sure. Alternatively, maybe I can suggest a division where I take the cups and they take the other items, which might align with their values.} \\
\\ 
\texttt{Looking at the items: there are 2 of each. If I take both cups, the remaining items are 2 wires, 2 glasses, 2 needles. The other player could take all of those, but since each item has only two copies, they can't take more than two of any. But if their value function is such that they value, say, wires and glasses, then taking those would be good for them. But I don't know.} \\
\\ 
\texttt{Another angle: maybe the other player's value function is such that they don't want cups. For example, if they value wires or glasses. If I take the cups, and they take the wires and glasses, that could work. But how to communicate that? I need to make a proposal that's likely to be complementary.} \\
\\ 
\texttt{Alternatively, maybe I should propose taking one cup and some other items, but that would lower my score. But if the other player also wants cups, maybe splitting the cups would allow both proposals to be complementary. For example, if I take one cup and they take the other, but then I get 4 points and they get whatever their value is. But I don't know their values. However, since my goal is to maximize my own score, taking both cups is better if possible.} \\
\\ 
\texttt{But the risk is that if they also want both cups, we both get zero. So maybe I should propose taking both cups and offer to give up one in exchange for something else. But how to do that without knowing their preferences.} \\
\\ 
\texttt{Wait, the rules say that if the proposals are complementary, meaning there's enough} \\
            }
        }
    }
    & & \\ \\

    \theutterance \stepcounter{utterance}  
    & & \multicolumn{4}{p{0.6\linewidth}}{
        \cellcolor[rgb]{0.9,0.9,0.9}{
            \makecell[{{p{\linewidth}}}]{
                \texttt{\tiny{[GM$\rangle$P2]}}
                \texttt{<think>} \\
\texttt{Okay, let's try to figure out the best way to approach this negotiation. My value function is cups at 4 each, wires at 1 each, and the others (glass and needle) at 0. So my top priority is to get as many cups as possible. The other player's values are secret, so I need to consider possible scenarios.} \\
\\ 
\texttt{First, if I take both cups, that gives me 8 points. But if the other player also wants cups, then our proposals would conflict, and we both get zero. So I need to find a way to coordinate. Maybe I can propose taking the cups and offer them other items in exchange. But since I don't know their preferences, I can't be sure. Alternatively, maybe I can suggest a division where I take the cups and they take the other items, which might align with their values.} \\
\\ 
\texttt{Looking at the items: there are 2 of each. If I take both cups, the remaining items are 2 wires, 2 glasses, 2 needles. The other player could take all of those, but since each item has only two copies, they can't take more than two of any. But if their value function is such that they value, say, wires and glasses, then taking those would be good for them. But I don't know.} \\
\\ 
\texttt{Another angle: maybe the other player's value function is such that they don't want cups. For example, if they value wires or glasses. If I take the cups, and they take the wires and glasses, that could work. But how to communicate that? I need to make a proposal that's likely to be complementary.} \\
\\ 
\texttt{Alternatively, maybe I should propose taking one cup and some other items, but that would lower my score. But if the other player also wants cups, maybe splitting the cups would allow both proposals to be complementary. For example, if I take one cup and they take the other, but then I get 4 points and they get whatever their value is. But I don't know their values. However, since my goal is to maximize my own score, taking both cups is better if possible.} \\
\\ 
\texttt{But the risk is that if they also want both cups, we both get zero. So maybe I should propose taking both cups and offer to give up one in exchange for something else. But how to do that without knowing their preferences.} \\
\\ 
\texttt{Wait, the rules say that if the proposals are complementary, meaning there's enough} \\
            }
        }
    }
    & \\ \\

    \theutterance \stepcounter{utterance}  
    & & & \multicolumn{4}{p{0.6\linewidth}}{
        \cellcolor[rgb]{0.9,0.9,0.9}{
            \makecell[{{p{\linewidth}}}]{
                \texttt{\tiny{[GM$\langle$P2]}}
                \texttt{<think>} \\
\texttt{Okay, let's try to figure out the best way to approach this negotiation. My value function is cups at 4 each, wires at 1 each, and the others (glass and needle) at 0. So my top priority is to get as many cups as possible. The other player's values are secret, so I need to consider possible scenarios.} \\
\\ 
\texttt{First, if I take both cups, that gives me 8 points. But if the other player also wants cups, then our proposals would conflict, and we both get zero. So I need to find a way to coordinate. Maybe I can propose taking the cups and offer them other items in exchange. But since I don't know their preferences, I can't be sure. Alternatively, maybe I can suggest a division where I take the cups and they take the other items, which might align with their values.} \\
\\ 
\texttt{Looking at the items: there are 2 of each. If I take both cups, the remaining items are 2 wires, 2 glasses, 2 needles. The other player could take all of those, but since each item has only two copies, they can't take more than two of any. But if their value function is such that they value, say, wires and glasses, then taking those would be good for them. But I don't know.} \\
\\ 
\texttt{Another angle: maybe the other player's value function is such that they don't want cups. For example, if they value wires or glasses. If I take the cups, and they take the wires and glasses, that could work. But how to communicate that? I need to make a proposal that's likely to be complementary.} \\
\\ 
\texttt{Alternatively, maybe I should propose taking one cup and some other items, but that would lower my score. But if the other player also wants cups, maybe splitting the cups would allow both proposals to be complementary. For example, if I take one cup and they take the other, but then I get 4 points and they get whatever their value is. But I don't know their values. However, since my goal is to maximize my own score, taking both cups is better if possible.} \\
\\ 
\texttt{But the risk is that if they also want both cups, we both get zero. So maybe I should propose taking both cups and offer to give up one in exchange for something else. But how to do that without knowing their preferences.} \\
\\ 
\texttt{Wait, the rules say that if the proposals are complementary, meaning there's enough} \\
            }
        }
    }
     \\ \\

    \theutterance \stepcounter{utterance}  
    & & & \multicolumn{2}{p{0.3\linewidth}}{
        \cellcolor[rgb]{0.9,0.9,0.9}{
            \makecell[{{p{\linewidth}}}]{
                \texttt{\tiny{[GM$|$GM]}}
                \texttt{continue} \\
            }
        }
    }
    & & \\ \\

    \theutterance \stepcounter{utterance}  
    & & & \multicolumn{2}{p{0.3\linewidth}}{
        \cellcolor[rgb]{0.9,0.9,0.9}{
            \makecell[{{p{\linewidth}}}]{
                \texttt{\tiny{[GM$|$GM]}}
                \texttt{<think>} \\
\texttt{Okay, let's try to figure out the best way to approach this negotiation. My value function is cups at 4 each, wires at 1 each, and the others (glass and needle) at 0. So my top priority is to get as many cups as possible. The other player's values are secret, so I need to consider possible scenarios.} \\
\\ 
\texttt{First, if I take both cups, that gives me 8 points. But if the other player also wants cups, then our proposals would conflict, and we both get zero. So I need to find a way to coordinate. Maybe I can propose taking the cups and offer them other items in exchange. But since I don't know their preferences, I can't be sure. Alternatively, maybe I can suggest a division where I take the cups and they take the other items, which might align with their values.} \\
\\ 
\texttt{Looking at the items: there are 2 of each. If I take both cups, the remaining items are 2 wires, 2 glasses, 2 needles. The other player could take all of those, but since each item has only two copies, they can't take more than two of any. But if their value function is such that they value, say, wires and glasses, then taking those would be good for them. But I don't know.} \\
\\ 
\texttt{Another angle: maybe the other player's value function is such that they don't want cups. For example, if they value wires or glasses. If I take the cups, and they take the wires and glasses, that could work. But how to communicate that? I need to make a proposal that's likely to be complementary.} \\
\\ 
\texttt{Alternatively, maybe I should propose taking one cup and some other items, but that would lower my score. But if the other player also wants cups, maybe splitting the cups would allow both proposals to be complementary. For example, if I take one cup and they take the other, but then I get 4 points and they get whatever their value is. But I don't know their values. However, since my goal is to maximize my own score, taking both cups is better if possible.} \\
\\ 
\texttt{But the risk is that if they also want both cups, we both get zero. So maybe I should propose taking both cups and offer to give up one in exchange for something else. But how to do that without knowing their preferences.} \\
\\ 
\texttt{Wait, the rules say that if the proposals are complementary, meaning there's enough} \\
            }
        }
    }
    & & \\ \\

    \theutterance \stepcounter{utterance}  
    & & \multicolumn{4}{p{0.6\linewidth}}{
        \cellcolor[rgb]{0.9,0.9,0.9}{
            \makecell[{{p{\linewidth}}}]{
                \texttt{\tiny{[P1$\langle$GM]}}
                \texttt{<think>} \\
\texttt{Okay, let's try to figure out the best way to approach this negotiation. My value function is cups at 4 each, wires at 1 each, and the others (glass and needle) at 0. So my top priority is to get as many cups as possible. The other player's values are secret, so I need to consider possible scenarios.} \\
\\ 
\texttt{First, if I take both cups, that gives me 8 points. But if the other player also wants cups, then our proposals would conflict, and we both get zero. So I need to find a way to coordinate. Maybe I can propose taking the cups and offer them other items in exchange. But since I don't know their preferences, I can't be sure. Alternatively, maybe I can suggest a division where I take the cups and they take the other items, which might align with their values.} \\
\\ 
\texttt{Looking at the items: there are 2 of each. If I take both cups, the remaining items are 2 wires, 2 glasses, 2 needles. The other player could take all of those, but since each item has only two copies, they can't take more than two of any. But if their value function is such that they value, say, wires and glasses, then taking those would be good for them. But I don't know.} \\
\\ 
\texttt{Another angle: maybe the other player's value function is such that they don't want cups. For example, if they value wires or glasses. If I take the cups, and they take the wires and glasses, that could work. But how to communicate that? I need to make a proposal that's likely to be complementary.} \\
\\ 
\texttt{Alternatively, maybe I should propose taking one cup and some other items, but that would lower my score. But if the other player also wants cups, maybe splitting the cups would allow both proposals to be complementary. For example, if I take one cup and they take the other, but then I get 4 points and they get whatever their value is. But I don't know their values. However, since my goal is to maximize my own score, taking both cups is better if possible.} \\
\\ 
\texttt{But the risk is that if they also want both cups, we both get zero. So maybe I should propose taking both cups and offer to give up one in exchange for something else. But how to do that without knowing their preferences.} \\
\\ 
\texttt{Wait, the rules say that if the proposals are complementary, meaning there's enough} \\
            }
        }
    }
    & \\ \\

    \theutterance \stepcounter{utterance}  
    & \multicolumn{4}{p{0.6\linewidth}}{
        \cellcolor[rgb]{0.9,0.9,0.9}{
            \makecell[{{p{\linewidth}}}]{
                \texttt{\tiny{[P1$\rangle$GM]}}
                \texttt{<think>} \\
\texttt{Okay, let's try to figure out the best way to approach this negotiation. My value function is cups at 4 each, wires at 1 each, and the others (glass and needle) at 0. So my top priority is to get as many cups as possible. The other player's values are secret, so I need to consider possible scenarios.} \\
\\ 
\texttt{First, if I take both cups, that gives me 8 points. But if the other player also wants cups, then our proposals would conflict, and we both get zero. So I need to find a way to coordinate. Maybe I can propose taking the cups and offer them other items in exchange. But since I don't know their preferences, I can't be sure. Alternatively, maybe I can suggest a division where I take the cups and they take the other items, which might align with their values.} \\
\\ 
\texttt{Looking at the items: there are 2 of each. If I take both cups, the remaining items are 2 wires, 2 glasses, 2 needles. The other player could take all of those, but since each item has only two copies, they can't take more than two of any. But if their value function is such that they value, say, wires and glasses, then taking those would be good for them. But I don't know.} \\
\\ 
\texttt{Another angle: maybe the other player's value function is such that they don't want cups. For example, if they value wires or glasses. If I take the cups, and they take the wires and glasses, that could work. But how to communicate that? I need to make a proposal that's likely to be complementary.} \\
\\ 
\texttt{Alternatively, maybe I should propose taking one cup and some other items, but that would lower my score. But if the other player also wants cups, maybe splitting the cups would allow both proposals to be complementary. For example, if I take one cup and they take the other, but then I get 4 points and they get whatever their value is. But I don't know their values. However, since my goal is to maximize my own score, taking both cups is better if possible.} \\
\\ 
\texttt{But the risk is that if they also want both cups, we both get zero. So maybe I should propose taking both cups and offer to give up one in exchange for something else. But how to do that without knowing their preferences.} \\
\\ 
\texttt{Wait, the rules say that if the proposals are complementary, meaning there's enough} \\
            }
        }
    }
    & & \\ \\

    \theutterance \stepcounter{utterance}  
    & & & \multicolumn{2}{p{0.3\linewidth}}{
        \cellcolor[rgb]{0.9,0.9,0.9}{
            \makecell[{{p{\linewidth}}}]{
                \texttt{\tiny{[GM$|$GM]}}
                \texttt{continue} \\
            }
        }
    }
    & & \\ \\

    \theutterance \stepcounter{utterance}  
    & & & \multicolumn{2}{p{0.3\linewidth}}{
        \cellcolor[rgb]{0.9,0.9,0.9}{
            \makecell[{{p{\linewidth}}}]{
                \texttt{\tiny{[GM$|$GM]}}
                \texttt{<think>} \\
\texttt{Okay, let's try to figure out the best way to approach this negotiation. My value function is cups at 4 each, wires at 1 each, and the others (glass and needle) at 0. So my top priority is to get as many cups as possible. The other player's values are secret, so I need to consider possible scenarios.} \\
\\ 
\texttt{First, if I take both cups, that gives me 8 points. But if the other player also wants cups, then our proposals would conflict, and we both get zero. So I need to find a way to coordinate. Maybe I can propose taking the cups and offer them other items in exchange. But since I don't know their preferences, I can't be sure. Alternatively, maybe I can suggest a division where I take the cups and they take the other items, which might align with their values.} \\
\\ 
\texttt{Looking at the items: there are 2 of each. If I take both cups, the remaining items are 2 wires, 2 glasses, 2 needles. The other player could take all of those, but since each item has only two copies, they can't take more than two of any. But if their value function is such that they value, say, wires and glasses, then taking those would be good for them. But I don't know.} \\
\\ 
\texttt{Another angle: maybe the other player's value function is such that they don't want cups. For example, if they value wires or glasses. If I take the cups, and they take the wires and glasses, that could work. But how to communicate that? I need to make a proposal that's likely to be complementary.} \\
\\ 
\texttt{Alternatively, maybe I should propose taking one cup and some other items, but that would lower my score. But if the other player also wants cups, maybe splitting the cups would allow both proposals to be complementary. For example, if I take one cup and they take the other, but then I get 4 points and they get whatever their value is. But I don't know their values. However, since my goal is to maximize my own score, taking both cups is better if possible.} \\
\\ 
\texttt{But the risk is that if they also want both cups, we both get zero. So maybe I should propose taking both cups and offer to give up one in exchange for something else. But how to do that without knowing their preferences.} \\
\\ 
\texttt{Wait, the rules say that if the proposals are complementary, meaning there's enough} \\
            }
        }
    }
    & & \\ \\

    \theutterance \stepcounter{utterance}  
    & & \multicolumn{4}{p{0.6\linewidth}}{
        \cellcolor[rgb]{0.9,0.9,0.9}{
            \makecell[{{p{\linewidth}}}]{
                \texttt{\tiny{[GM$\rangle$P2]}}
                \texttt{<think>} \\
\texttt{Okay, let's try to figure out the best way to approach this negotiation. My value function is cups at 4 each, wires at 1 each, and the others (glass and needle) at 0. So my top priority is to get as many cups as possible. The other player's values are secret, so I need to consider possible scenarios.} \\
\\ 
\texttt{First, if I take both cups, that gives me 8 points. But if the other player also wants cups, then our proposals would conflict, and we both get zero. So I need to find a way to coordinate. Maybe I can propose taking the cups and offer them other items in exchange. But since I don't know their preferences, I can't be sure. Alternatively, maybe I can suggest a division where I take the cups and they take the other items, which might align with their values.} \\
\\ 
\texttt{Looking at the items: there are 2 of each. If I take both cups, the remaining items are 2 wires, 2 glasses, 2 needles. The other player could take all of those, but since each item has only two copies, they can't take more than two of any. But if their value function is such that they value, say, wires and glasses, then taking those would be good for them. But I don't know.} \\
\\ 
\texttt{Another angle: maybe the other player's value function is such that they don't want cups. For example, if they value wires or glasses. If I take the cups, and they take the wires and glasses, that could work. But how to communicate that? I need to make a proposal that's likely to be complementary.} \\
\\ 
\texttt{Alternatively, maybe I should propose taking one cup and some other items, but that would lower my score. But if the other player also wants cups, maybe splitting the cups would allow both proposals to be complementary. For example, if I take one cup and they take the other, but then I get 4 points and they get whatever their value is. But I don't know their values. However, since my goal is to maximize my own score, taking both cups is better if possible.} \\
\\ 
\texttt{But the risk is that if they also want both cups, we both get zero. So maybe I should propose taking both cups and offer to give up one in exchange for something else. But how to do that without knowing their preferences.} \\
\\ 
\texttt{Wait, the rules say that if the proposals are complementary, meaning there's enough} \\
            }
        }
    }
    & \\ \\

    \theutterance \stepcounter{utterance}  
    & & & \multicolumn{4}{p{0.6\linewidth}}{
        \cellcolor[rgb]{0.9,0.9,0.9}{
            \makecell[{{p{\linewidth}}}]{
                \texttt{\tiny{[GM$\langle$P2]}}
                \texttt{<think>} \\
\texttt{Okay, let's try to figure out the best way to approach this negotiation. My value function is cups at 4 each, wires at 1 each, and the others (glass and needle) at 0. So my top priority is to get as many cups as possible. The other player's values are secret, so I need to consider possible scenarios.} \\
\\ 
\texttt{First, if I take both cups, that gives me 8 points. But if the other player also wants cups, then our proposals would conflict, and we both get zero. So I need to find a way to coordinate. Maybe I can propose taking the cups and offer them other items in exchange. But since I don't know their preferences, I can't be sure. Alternatively, maybe I can suggest a division where I take the cups and they take the other items, which might align with their values.} \\
\\ 
\texttt{Looking at the items: there are 2 of each. If I take both cups, the remaining items are 2 wires, 2 glasses, 2 needles. The other player could take all of those, but since each item has only two copies, they can't take more than two of any. But if their value function is such that they value, say, wires and glasses, then taking those would be good for them. But I don't know.} \\
\\ 
\texttt{Another angle: maybe the other player's value function is such that they don't want cups. For example, if they value wires or glasses. If I take the cups, and they take the wires and glasses, that could work. But how to communicate that? I need to make a proposal that's likely to be complementary.} \\
\\ 
\texttt{Alternatively, maybe I should propose taking one cup and some other items, but that would lower my score. But if the other player also wants cups, maybe splitting the cups would allow both proposals to be complementary. For example, if I take one cup and they take the other, but then I get 4 points and they get whatever their value is. But I don't know their values. However, since my goal is to maximize my own score, taking both cups is better if possible.} \\
\\ 
\texttt{But the risk is that if they also want both cups, we both get zero. So maybe I should propose taking both cups and offer to give up one in exchange for something else. But how to do that without knowing their preferences.} \\
\\ 
\texttt{Wait, the rules say that if the proposals are complementary, meaning there's enough} \\
            }
        }
    }
     \\ \\

    \theutterance \stepcounter{utterance}  
    & & & \multicolumn{2}{p{0.3\linewidth}}{
        \cellcolor[rgb]{0.9,0.9,0.9}{
            \makecell[{{p{\linewidth}}}]{
                \texttt{\tiny{[GM$|$GM]}}
                \texttt{continue} \\
            }
        }
    }
    & & \\ \\

    \theutterance \stepcounter{utterance}  
    & & & \multicolumn{2}{p{0.3\linewidth}}{
        \cellcolor[rgb]{0.9,0.9,0.9}{
            \makecell[{{p{\linewidth}}}]{
                \texttt{\tiny{[GM$|$GM]}}
                \texttt{<think>} \\
\texttt{Okay, let's try to figure out the best way to approach this negotiation. My value function is cups at 4 each, wires at 1 each, and the others (glass and needle) at 0. So my top priority is to get as many cups as possible. The other player's values are secret, so I need to consider possible scenarios.} \\
\\ 
\texttt{First, if I take both cups, that gives me 8 points. But if the other player also wants cups, then our proposals would conflict, and we both get zero. So I need to find a way to coordinate. Maybe I can propose taking the cups and offer them other items in exchange. But since I don't know their preferences, I can't be sure. Alternatively, maybe I can suggest a division where I take the cups and they take the other items, which might align with their values.} \\
\\ 
\texttt{Looking at the items: there are 2 of each. If I take both cups, the remaining items are 2 wires, 2 glasses, 2 needles. The other player could take all of those, but since each item has only two copies, they can't take more than two of any. But if their value function is such that they value, say, wires and glasses, then taking those would be good for them. But I don't know.} \\
\\ 
\texttt{Another angle: maybe the other player's value function is such that they don't want cups. For example, if they value wires or glasses. If I take the cups, and they take the wires and glasses, that could work. But how to communicate that? I need to make a proposal that's likely to be complementary.} \\
\\ 
\texttt{Alternatively, maybe I should propose taking one cup and some other items, but that would lower my score. But if the other player also wants cups, maybe splitting the cups would allow both proposals to be complementary. For example, if I take one cup and they take the other, but then I get 4 points and they get whatever their value is. But I don't know their values. However, since my goal is to maximize my own score, taking both cups is better if possible.} \\
\\ 
\texttt{But the risk is that if they also want both cups, we both get zero. So maybe I should propose taking both cups and offer to give up one in exchange for something else. But how to do that without knowing their preferences.} \\
\\ 
\texttt{Wait, the rules say that if the proposals are complementary, meaning there's enough} \\
            }
        }
    }
    & & \\ \\

    \theutterance \stepcounter{utterance}  
    & & \multicolumn{4}{p{0.6\linewidth}}{
        \cellcolor[rgb]{0.9,0.9,0.9}{
            \makecell[{{p{\linewidth}}}]{
                \texttt{\tiny{[P1$\langle$GM]}}
                \texttt{<think>} \\
\texttt{Okay, let's try to figure out the best way to approach this negotiation. My value function is cups at 4 each, wires at 1 each, and the others (glass and needle) at 0. So my top priority is to get as many cups as possible. The other player's values are secret, so I need to consider possible scenarios.} \\
\\ 
\texttt{First, if I take both cups, that gives me 8 points. But if the other player also wants cups, then our proposals would conflict, and we both get zero. So I need to find a way to coordinate. Maybe I can propose taking the cups and offer them other items in exchange. But since I don't know their preferences, I can't be sure. Alternatively, maybe I can suggest a division where I take the cups and they take the other items, which might align with their values.} \\
\\ 
\texttt{Looking at the items: there are 2 of each. If I take both cups, the remaining items are 2 wires, 2 glasses, 2 needles. The other player could take all of those, but since each item has only two copies, they can't take more than two of any. But if their value function is such that they value, say, wires and glasses, then taking those would be good for them. But I don't know.} \\
\\ 
\texttt{Another angle: maybe the other player's value function is such that they don't want cups. For example, if they value wires or glasses. If I take the cups, and they take the wires and glasses, that could work. But how to communicate that? I need to make a proposal that's likely to be complementary.} \\
\\ 
\texttt{Alternatively, maybe I should propose taking one cup and some other items, but that would lower my score. But if the other player also wants cups, maybe splitting the cups would allow both proposals to be complementary. For example, if I take one cup and they take the other, but then I get 4 points and they get whatever their value is. But I don't know their values. However, since my goal is to maximize my own score, taking both cups is better if possible.} \\
\\ 
\texttt{But the risk is that if they also want both cups, we both get zero. So maybe I should propose taking both cups and offer to give up one in exchange for something else. But how to do that without knowing their preferences.} \\
\\ 
\texttt{Wait, the rules say that if the proposals are complementary, meaning there's enough} \\
\\ 
\\ 
\texttt{The time is up. It is now your turn to submit a secret proposal.} \\
            }
        }
    }
    & \\ \\

    \theutterance \stepcounter{utterance}  
    & \multicolumn{4}{p{0.6\linewidth}}{
        \cellcolor[rgb]{0.9,0.9,0.9}{
            \makecell[{{p{\linewidth}}}]{
                \texttt{\tiny{[P1$\rangle$GM]}}
                \texttt{<think>} \\
\texttt{Okay, let's try to figure out the best way to approach this negotiation. My value function is cups at 4 each, wires at 1 each, and the others (glass and needle) at 0. So my top priority is to get as many cups as possible. The other player's values are secret, so I need to consider possible scenarios.} \\
\\ 
\texttt{First, if I take both cups, that gives me 8 points. But if the other player also wants cups, then our proposals would conflict, and we both get zero. So I need to find a way to coordinate. Maybe I can propose taking the cups and offer them other items in exchange. But since I don't know their preferences, I can't be sure. Alternatively, maybe I can suggest a division where I take the cups and they take the other items, which might align with their values.} \\
\\ 
\texttt{Looking at the items: there are 2 of each. If I take both cups, the remaining items are 2 wires, 2 glasses, 2 needles. The other player could take all of those, but since each item has only two copies, they can't take more than two of any. But if their value function is such that they value, say, wires and glasses, then taking those would be good for them. But I don't know.} \\
\\ 
\texttt{Another angle: maybe the other player's value function is such that they don't want cups. For example, if they value wires or glasses. If I take the cups, and they take the wires and glasses, that could work. But how to communicate that? I need to make a proposal that's likely to be complementary.} \\
\\ 
\texttt{Alternatively, maybe I should propose taking one cup and some other items, but that would lower my score. But if the other player also wants cups, maybe splitting the cups would allow both proposals to be complementary. For example, if I take one cup and they take the other, but then I get 4 points and they get whatever their value is. But I don't know their values. However, since my goal is to maximize my own score, taking both cups is better if possible.} \\
\\ 
\texttt{But the risk is that if they also want both cups, we both get zero. So maybe I should propose taking both cups and offer to give up one in exchange for something else. But how to do that without knowing their preferences.} \\
\\ 
\texttt{Wait, the rules say that if the proposals are complementary, meaning there's enough} \\
            }
        }
    }
    & & \\ \\

    \theutterance \stepcounter{utterance}  
    & & & \multicolumn{2}{p{0.3\linewidth}}{
        \cellcolor[rgb]{0.9,0.9,0.9}{
            \makecell[{{p{\linewidth}}}]{
                \texttt{\tiny{[GM$|$GM]}}
                \texttt{continue} \\
            }
        }
    }
    & & \\ \\

    \theutterance \stepcounter{utterance}  
    & & & \multicolumn{2}{p{0.3\linewidth}}{
        \cellcolor[rgb]{0.9,0.9,0.9}{
            \makecell[{{p{\linewidth}}}]{
                \texttt{\tiny{[GM$|$GM]}}
                \texttt{player was instructed to make a proposal} \\
            }
        }
    }
    & & \\ \\

\end{supertabular}
}

\end{document}
